According to variation theory \autocite[See][]{NCOL}, learning an educational
objective requires the learner to discern all critical aspects of that
objective. Misconceptions occur when the learner cannot yet discern one or
more critical aspects; documented misconceptions therefore help us identify
what those aspects are.

In this work (in progress), we apply this approach---developed for
introductory programming in our companion paper
\autocite{VTProgMisconceptions}---to the Unix shell: the command model, the
hierarchical file system, composition (pipes, redirection, scripts), and
remote shells (SSH). We review the literature on learners' difficulties and
mental models of the command line, treat the documented misconceptions as
phenomenographic observations, analyse them through the lens of variation
theory to identify critical aspects, and outline patterns of variation to
teach learners to discern these aspects. The shell is a particularly
interesting case: research on how students learn it is scarce even though it
is widely taught, and the command line's vary-and-observe loop makes it a
natural medium for variation-theory-based teaching.
